\section{The turning point and the missing quarter-period}\label{sec:turning}

This section follows the frame angle $\chi_m$ from site $0$ to a block $M$ well inside the
oscillatory region and shows that it equals the discrete Bohr--Sommerfeld sum up to $M$,
plus a drift, with no quarter-period (\Cref{prop:nearfield}). We first explain the
result as a budget of phase, region by region, and then prove the three entries of the
budget. The argument is the same for balanced and dimerised links, so we treat them
together under the following hypothesis.

\begin{hypothesis}\label{hyp:Hk}
$a_n=c\,n^\alpha\bigl(1+\kappa'(-1)^n/n+r_n\bigr)$ with $r$ as in \Cref{hyp:H} and
$2|\kappa'|<\alpha-1$; the diagonal satisfies $|b_n|\le C_b(1+n)^\beta$ with
$\beta<\alpha-1-2|\kappa'|$.
\end{hypothesis}

\noindent The odd links $A_m=a_{2m+1}$ see the alternating term only through a constant:
$A_m=c\,N_m^\alpha(1+\hat r_{2m+1})$ with $\hat r_n:=r_n-\kappa'/n$, and $\hat r$ again
satisfies \Cref{hyp:H}. We write $\bar\ell_m:=\alpha\log(N_{m+1}/N_m)$ for the increment of
$\log(c\,N^\alpha)$ and $\Delta\hat e_m:=\log(1+\hat r_{2m+3})-\log(1+\hat r_{2m+1})$,
so that $\log A_{m+1}-\log A_m=\bar\ell_m+\Delta\hat e_m$ with $|\Delta\hat e_m|\le CN_m^{-1-\delta}$.

\subsection{The phase budget}

Measure the distance from the turning point on the Airy scale
\begin{equation}\label{eq:ell}
  \ell:=\bigl(N_\lambda/(16\alpha)\bigr)^{1/3},\qquad x_m:=4\ell^2(\mu_m-2)=-4\ell^2u_m .
\end{equation}
Then $x_m>0$ on the \emph{hyperbolic} side, where $2A_m<\lambda$, and $x_m<0$ on the
\emph{oscillatory} side. Near the turning point $x_m$ decreases by $\ell^{-1}(1+o(1))$ per
block (\Cref{lem:airycoords}), so it is the variable in which the turning point has a fixed
shape, and the window $|x_m|<K$ is about $2K\ell$ blocks, that is $O(N_\lambda^{1/3})$ sites,
wide. Fix $K\ge1$, let $m_L$ be the largest block with $x_m\ge K$ and $m_R$ the least
with $x_m\le-K$. The blocks then fall into the regions sketched in \Cref{fig:budget}, and
each contributes a definite amount to $\chi$.

\emph{Below the turning point each block contributes $\pi$.} For $m\le m_L$ the regular
solution has no sign change: it is positive and increasing (\Cref{lem:hyperbolic}). The
reference solutions $\Pt_{2m}$ and $\Rt_{2m+1}$, on the other hand, change sign from each
block to the next, by \eqref{eq:PR}. In the coordinates of \Cref{thm:count}, which expand
$p$ in $\Pt$ and $\Rt$, the phase therefore advances by $\pi$ per block in the following sense: after
block $m$ the lifted angle $\psi_{2m-1}$ lies in the open quarter $(m\pi-\pi/2,m\pi)$. The increment of a
single block is not exactly $\pi$; only the quarter is fixed, and after the hyperbolic region
$\psi_{2m_L-1}$ lies in $(m_L\pi-\pi/2,m_L\pi)$. Since
$\varphi_m=\pi$ there too, the hyperbolic region contributes exactly its terms $m_L\pi$ of
the sum $S$, up to a position inside one quarter that the window decides.

\emph{The Airy window contributes $\Omega(K)-\pi/4$.} Across the window the block is a small
perturbation of $T_1(2)=-I+\mathcal N$, where $\mathcal N$ is nilpotent with kernel
$e_-=(1,-1)/\sqrt2$, and each block acts on the pair (amplitude, slope) as an Euler step
of length $\ell^{-1}$ for the Airy equation $y''=xy$ (\Cref{lem:airycoords}). The amplitude
$y$ is the component along $e_-$. The regular
solution enters the window as the recessive solution $\Ai$, because it was increasing,
and leaves it at $x=-K$ with the Airy phase
$\Omega(K)=\arg\bigl(\Ai(-K)-i\Ai'(-K)/\sqrt K\bigr)$, measured from the direction $e_-$.
That direction has angle $-\pi/4$, so (\Cref{blem:window})
\[
  \chi_{m_R}-m_R\pi\ \longrightarrow\ \Omega(K)-\frac\pi4 .
\]
In the same window the sum $S$ falls short of $\pi$ per block by
$T_K:=\sum_{m_L<m<m_R}(\pi-\varphi_m)\to\frac23K^{3/2}$ (\Cref{lem:turningsum}), which is the
classical action across the window. Hence
\begin{equation}\label{eq:budget-window}
  \chi_{m_R}-\sum_{m<m_R}\varphi_m\ \longrightarrow\ \Omega(K)-\frac\pi4+\frac23K^{3/2}
  =O(K^{-3/2}),
\end{equation}
because $\Omega(K)=\frac\pi4-\frac23K^{3/2}+O(K^{-3/2})$ (\Cref{lem:airyphase}). The action
$\frac23K^{3/2}$ is already in $S$, and what is left over is $\frac\pi4-\frac\pi4$.

\emph{The two quarter-periods.} The $+\frac\pi4$ in $\Omega$ is the Maslov phase of the
turning point, the $-\frac\pi4$ in $\Ai(-\zeta)\sim\pi^{-1/2}\zeta^{-1/4}\cos(\frac23\zeta^{3/2}-\frac\pi4)$.
The $-\frac\pi4$ is set by the limit-circle end. At the turning point $\lambda=2a$ the
local solution of $a u_{k-1}+au_{k+1}=\lambda u_k$ is the band-edge wave $u_k=\mathrm{const}$,
which takes equal values on the two sites of a block. The coordinates of
\Cref{thm:count} measure a solution against $\Pt$ and $\Rt$, which live on alternate
sites. The block frame $E_m=D_m\Et(\mu_m)$ first rescales them by $D_m$ to equal amplitude
on block $m$, and in that scaling a vector with equal values on the two sites has angle
$-\frac\pi4$ modulo $\pi$: it is the direction $e_-$. For balanced links the frames tend to
the identity as $m\to\infty$, and $D_m\to I$ precisely because $\Pt$ and $\Rt$ were normalised by $\vp$ to have
equal amplitude at infinity (\Cref{prop:zero}(c)). So the $-\frac\pi4$ measured at the
turning point is the $-\frac\pi4$ read at the limit-circle end, in the coordinates in which
the boundary condition of $H_t$ is imposed. It is the band-edge direction in the staggered
equal-amplitude coordinates, not an independent Maslov index of the end: the absence of an extra
constant is relative to this normalisation, which is fixed once by the data at $\lambda=0$. A change of
the boundary coordinates changes the angle of a given extension and its offset together, without
changing any spectrum. (For dimerised links $D_m$ keeps drifting
like $\operatorname{diag}(N_m^{\kappa'},N_m^{-\kappa'})$, and \Cref{sec:dimer} follows that drift
to the end.) The cancellation takes place in $\psi$, before an extension is chosen, which is why every
extension obeys a law $j\pi+t$ and none obeys $(j+\frac14)\pi+t$.

\emph{Beyond the window each block adds $\varphi_m$ and a drift.} On the oscillatory side
the adiabatic normal form (\Cref{prop:adiabatic}) applies. Each block adds its rotation
$\varphi_m$ and the frame drift $\bar\varepsilon_m\tan(\varphi_m/2)$ with
$\bar\varepsilon_m=-2\kappa'/N_m$, and the drifts add up to $-\kappa'\pi/(2\alpha)$. This is
$0$ for balanced links.

Adding the entries gives the near-field law
\begin{equation}\label{eq:budget}
\begin{aligned}
  \chi_M&=\underbrace{m_L\pi}_{\text{hyperbolic}}
  +\underbrace{\sum_{m_L\le m<m_R}\varphi_m+\Bigl(\Omega(K)-\frac\pi4+T_K\Bigr)}_{\text{Airy window}}
  +\underbrace{\sum_{m_R\le m<M}\varphi_m-\frac{\kappa'\pi}{2\alpha}}_{\text{oscillatory}}+o(1)\\
  &=\sum_{m<M}\varphi_m-\frac{\kappa'\pi}{2\alpha}+o(1),
\end{aligned}
\end{equation}
where $o(1)$ is taken as $\lambda\to\infty$ and then $K\to\infty$. For balanced links $M=\infty$
and $\chi_\infty=\psi$, so $\psi=S+o(1)$. For dimerised links the frames do not converge, and
the far field beyond $M$ is the subject of \Cref{sec:dimer}.

\begin{schemfigure}
\begin{tikzpicture}
\begin{axis}[es plot, width=0.97\linewidth, height=7.6cm,
  xmode=log, ymode=log, xmin=1, xmax=1e6, ymin=2e-7, ymax=60,
  xlabel={site $N$ (logarithmic)}, ylabel={angle per block},
  xtick={1,10,100,1000,10000,100000,1000000},
  ytick={1e-6,1e-4,1e-2,1}, ymajorgrids=false,
  clip=false]
  % regions
  \addplot[draw=none, fill=esPlotGrid, fill opacity=0.6] coordinates {(88,2e-7) (112,2e-7) (112,60) (88,60)} \closedcycle;
  \addplot[draw=none, fill=esPlotGrid, fill opacity=0.35] coordinates {(1e4,2e-7) (1e6,2e-7) (1e6,60) (1e4,60)} \closedcycle;
  % rotation per block 2 arcsin(min(1,(N_lambda/N)^alpha)), alpha = 2, N_lambda = 100;
  % beyond N = 200 the arcsine is evaluated by its series u + u^3/6 + 3u^5/40,
  % because pgfmath loses the arcsine of small arguments
  \addplot[esPlotBlue, line width=1.2pt, samples=300, domain=0:2.30103, variable=\t]
    ({10^\t},{2*rad(asin(min(1,(100/10^\t)^2)))});
  \addplot[esPlotBlue, line width=1.2pt, samples=200, domain=2.30103:5.45, variable=\t]
    ({10^\t},{2*(100/10^\t)^2*(1+(100/10^\t)^4/6+3*(100/10^\t)^8/40)});
  % frame drift 2|kappa'|/N, kappa' = 0.25
  \addplot[esPlotAmber, line width=1.1pt, dashed, domain=0:6, samples=20, variable=\t]
    ({10^\t},{0.5/10^\t});
  \draw[esPlotGrey, densely dotted] (axis cs:100,2e-7) -- (axis cs:100,60);
  \draw[esPlotGrey, densely dotted] (axis cs:1e4,2e-7) -- (axis cs:1e4,60);
  \node[font=\scriptsize, text=esPlotInk, anchor=south] at (axis cs:100,60) {$N_\lambda$};
  \node[font=\scriptsize, text=esPlotInk, anchor=south] at (axis cs:1e4,60) {$N_{\rm dec}$ ($\tau=1$)};
  % labels on the curves
  \node[font=\scriptsize, text=esPlotBlue, anchor=south west] at (axis cs:1000,1.2e-2) {rotation $\varphi_m$};
  \node[font=\scriptsize, text=esPlotAmber, anchor=north west] at (axis cs:1.6,2.5e-3) {frame drift $|\varepsilon_m|\approx2|\kappa'|/N_m$};
  % budget annotations
  \node[font=\scriptsize, align=center, text=esPlotInk, fill=esPlotPanel, inner sep=2pt] at (axis cs:7,0.15)
    {hyperbolic\\ $\varphi_m=\pi$; $\psi$ advances\\ by $\pi$ per block};
  \node[font=\scriptsize, align=left, text=esPlotInk, fill=esPlotPanel, inner sep=2pt, anchor=west] at (axis cs:125,14)
    {Airy window, $\asymp N_\lambda^{1/3}$ sites:\\ $\Omega(K)-\frac\pi4+T_K\to0$};
  \node[font=\scriptsize, align=center, text=esPlotInk, fill=esPlotPanel, inner sep=2pt] at (axis cs:1100,3e-6)
    {near field: $\varphi_m$ plus drift,\\ drift total $-\kappa'\pi/(2\alpha)$};
  \node[font=\scriptsize, align=center, text=esPlotInk, fill=esPlotPanel, inner sep=2pt] at (axis cs:1e5,2e-2)
    {far field ($\kappa'\ne0$):\\ Bessel of order $\nu$};
\end{axis}
\end{tikzpicture}
\end{schemfigure}
\esfiglabel{fig:budget}%
\schemcap{FIG.\ \theesfig}{The phase budget along the chain, drawn for $\alpha=2$, $N_\lambda=100$ and
$\kappa'=\frac14$. Below the turning point $N_\lambda$ each block contributes $\pi$ to the sum $S$, and
the phase $\psi$ advances by $\pi$ per block up to a position inside one quarter. In the Airy
window the Maslov phase $\frac\pi4$ of the turning point cancels the $-\frac\pi4$ fixed by
the limit-circle end in the equal-amplitude coordinates set at $\lambda=0$, so no quarter-period
survives in those coordinates. Beyond it each block adds its
rotation and a drift. Where the rotation falls below the frame drift, near $N_{\rm dec}$, a
dimerised chain becomes a Bessel problem (\Cref{sec:dimer}).}

The rest of the section proves the three entries. The block geometry and the error sums
that the proofs need are collected in \Cref{app:routine}.

\subsection{The hyperbolic side}

\begin{lemma}[Monotone region]\label{lem:hyperbolic}
If $\lambda-b_k\ge a_k+a_{k+1}$ for $0\le k\le k_1$, then
$1=p_0\le p_1\le\dots\le p_{k_1+1}$ and $\psi_k\in(k\pi/2,(k+1)\pi/2)$ for $0\le k\le k_1$.
\end{lemma}

\begin{proof}
Inductively,
\[
  p_{k+1}=\frac{(\lambda-b_k)p_k-a_kp_{k-1}}{a_{k+1}}\ge\frac{(\lambda-b_k-a_k)\,p_k}{a_{k+1}}\ge p_k .
\]
By \Cref{thm:count}(i), $v_k$ is a positive multiple of
\[
  \Bigl(\frac{p_{2m}}{\Pt_{2m}},\frac{p_{2m+1}}{\Rt_{2m+1}}\Bigr)\quad(k=2m),\qquad
  \Bigl(\frac{p_{2m+2}}{\Pt_{2m+2}},\frac{p_{2m+1}}{\Rt_{2m+1}}\Bigr)\quad(k=2m+1),
\]
whose entries have the signs $(-1)^m,(-1)^m$ and $(-1)^{m+1},(-1)^m$ respectively, by \eqref{eq:PR}. In
both cases the angle of $v_k$ lies in $(k\pi/2,(k+1)\pi/2)$ modulo $2\pi$. All
$\gamma_k>0$, so each canonical increment lies in $[0,\pi)$; starting from
$\psi_0\in(0,\pi/2)$, the lift $\psi_k$ lies in an interval of length $3\pi/2$ that meets
the class $(k\pi/2,(k+1)\pi/2)+2\pi\ZZ$ only once.
\end{proof}

\noindent For $\lambda\ge\lambda_0(K)$ the hypothesis holds with $k_1=2m_L$, and therefore
$\psi_{2m_L-1}\in(m_L\pi-\pi/2,m_L\pi)$ (\Cref{lem:startwindow}).

\subsection{The Airy window}

\begin{boxedlemma}{Airy window}{window}
Assume \Cref{hyp:Hk} and fix $K\ge1$. Then
\[
  \limsup_{\lambda\to\infty}\Bigl|\chi_{m_R}-m_R\pi-\bigl(\Omega(K)-\tfrac\pi4\bigr)\Bigr|
  \le C_0\,\eta_A(K),\qquad \eta_A(K):=(1+\sqrt K)\,\frac{\Ai(K)}{\Bi(K)},
\]
where $\Omega$ is the Airy phase of \Cref{lem:airyphase} and $C_0$ is absolute.
\end{boxedlemma}

The proof works in coordinates adapted to the turning point, in which the blocks become
Euler steps of the Airy equation. Write $z_m:=D_m^{-1}v_{2m-1}$ and define $(X_m,Y_m)$ by
\begin{equation}\label{eq:XY}
  z_m=(-1)^m\bigl(X_m\,e_-+\tfrac{Y_m}{4\ell}\,e_+\bigr),\qquad\text{that is}\qquad
  (X_m,Y_m)^{\mathsf T}=(-1)^mC_\ell^{-1}z_m,\quad C_\ell:=\bigl[e_-,\ e_+/(4\ell)\bigr].
\end{equation}
The pair $(X_m,Y_m)$ will track $(y(x_m),-y'(x_m))$ for a solution $y$ of the Airy
equation $y''=xy$. A direct computation of $D_m^{-1}v_{2m-1}$ from \Cref{thm:count}(i)
and \eqref{eq:PR} gives the exact form
\begin{equation}\label{eq:zm}
  z_m\ \propto\ (-1)^mA_m^{1/2}\bigl(p_{2m},\,-c'_m\,p_{2m-1}\bigr),\qquad
  c'_m:=a_{2m}/a_{2m+1},
\end{equation}
with a positive constant of proportionality; no normalisation by $\vp$ or $\rho_m$
remains. So $X_m$ is, up to a factor, the average of $p$ over the block and $Y_m/\ell$ its
difference. The \emph{Airy angle} of block $m$ is
$\varrho_m:=\hl C_\ell^{-1}\bigl(\hl D_m^{-1}(\psi_{2m-1})-m\pi\bigr)$, where $\hl C_\ell$ is the
lift of $C_\ell$ with $\hl C_\ell(0)=-\pi/4$, the angle of $e_-$. \Cref{lem:airycoords}
shows that the blocks act on $(X,Y)$ as Euler steps $I+h_m\mathsf B(x_m)$ of the Airy system, with
$h_m=x_m-x_{m+1}=\ell^{-1}(1+o(1))$ and $\mathsf B(x)=\bigl(\begin{smallmatrix}0&1\\x&0\end{smallmatrix}\bigr)$,
that $\varrho_{m+1}$ is the principal lift of one step applied to $\varrho_m$, and that
the window starts with $X_{m_L}>0$ and $Y_{m_L}/X_{m_L}\ge-C\ell/N_\lambda$.

\begin{proof}[Proof of \Cref{blem:window}]
The argument has three parts: the discrete trajectory follows an Airy trajectory, that
trajectory is almost exactly $\Ai$, and its phase is read in the frame at $m_R$.

\step{the Airy flow} Let $\mathcal U(x,x')$ be the propagator of $(X,Y)'=-(Y,xX)$ in the
decreasing variable $x$, which is the system satisfied by $(y,-y')$ for $y''=xy$. Then
$\|\mathcal U(x_{m+1},x_m)-I-h_m\mathsf B(x_m)\|\le C(K)h_m^2$. With \Cref{lem:airycoords}(a),
telescoping gives
$\|\prod_{m'<m}M_{m'}-\mathcal U(x_m,x_{m_L})\|\le e^{C(K)}\bigl(C(K)\ell^{-1}+o(1)\bigr)\to0$
uniformly in $m$ and in the initial direction. The discrete vectors therefore stay
uniformly close to an Airy trajectory and away from $0$, and the principal steps
$\varrho_m$ stay within $o(1)$ of the continuous lifted angle of that trajectory.

\step{the trajectory is $\Ai$} Every limiting trajectory starts at $x=K$ with $X>0$ and
$Y/X=:\varsigma\ge0$, so it is a multiple of $y=\Ai+\epsilon\Bi$ with
$\epsilon=(\varsigma\Ai(K)+\Ai'(K))/(-\Bi'(K)-\varsigma\Bi(K))$. Since $\Bi,\Bi'>0$, the mediant
inequality gives $|\epsilon|\le\max(\Ai(K)/\Bi(K),|\Ai'(K)|/\Bi'(K))\le C_0\eta_A(K)$ for
every $\varsigma\in[0,\infty]$ (\cite{DLMF}, 9.7.5--9.7.8; equation numbers of the form 9.7.5
refer to that source throughout). The Wronskian $\Ai\Bi'-\Ai'\Bi=1/\pi$ gives
$\Ai(K)+\epsilon\Bi(K)=\pi^{-1}/(\Bi'(K)+\varsigma\Bi(K))>0$, and since $\Ai/\Bi$ decreases on
$[0,\infty)$, $y>0$ on $[0,K]$. So $\varrho$ stays in $(-\pi/2,\pi/2)$ on the hyperbolic part
of the window.

\step{reading the phase at $m_R$} By definition $\chi_{m_R}=\hl{\Et}(\mu_{m_R})^{-1}(\hl D_{m_R}^{-1}(\psi_{2m_R-1}))$,
so $\chi_{m_R}-m_R\pi$ is the image of $\varrho_{m_R}$ under the lift of $\Et(\mu_{m_R})^{-1}C_\ell$
that sends $0$ to $-\pi/4$. Since $\Et^{-1}$ multiplies $e_-$ by $e^{-s}$ and $e_+$ by $e^{s}$,
this map sends $(X,Y)$ to a positive multiple of $Xe_-+qYe_+$ with $q=e^{2s}/(4\ell)$, and
$q=\sqrt{4-u}/(2\sqrt{|x_{m_R}|})\to1/\sqrt K$. Hence
$\chi_{m_R}-m_R\pi\to-\pi/4+\Omega_\epsilon(K)$, where $\Omega_\epsilon(\zeta)$ is the
continuous lift of $\arg\bigl(y(-\zeta)-iy'(-\zeta)/\sqrt\zeta\bigr)$ on $(0,K]$, equal to
$\pi/2$ at $0^+$. The Wronskian and DLMF 9.7.9, 9.7.11 give
$|(y,y'/\sqrt K)(-K)|\ge cK^{-1/4}$, so $|\Omega_\epsilon(K)-\Omega_0(K)|\le C|\epsilon|$.
\end{proof}

\begin{lemma}[Airy phase]\label{lem:airyphase}
Let $\Omega$ be the continuous lift on $(0,\infty)$ of
$\arg\bigl(\Ai(-\zeta)-i\Ai'(-\zeta)/\sqrt\zeta\bigr)$ with $\Omega(0^+)=\pi/2$. Then
$\Omega(\zeta)=\pi/4-\tfrac23\zeta^{3/2}+O(\zeta^{-3/2})$.
\end{lemma}

\begin{proof}
By DLMF 9.7.9 and 9.7.10, $\Ai(-\zeta)=\pi^{-1/2}\zeta^{-1/4}(\cos\phi+O(\zeta^{-3/2}))$ and
$\Ai'(-\zeta)=\pi^{-1/2}\zeta^{1/4}(\sin\phi+O(\zeta^{-3/2}))$ with $\phi=\frac23\zeta^{3/2}-\frac\pi4$.
Hence
\[
  \Ai(-\zeta)-i\Ai'(-\zeta)/\sqrt\zeta=\pi^{-1/2}\zeta^{-1/4}\bigl(e^{-i\phi}+O(\zeta^{-3/2})\bigr),
\]
and $\Omega(\zeta)=-\phi+2\pi n_0+O(\zeta^{-3/2})$ for a fixed integer $n_0$. Writing
$\Ai(-\zeta)=r\cos\Omega$, $-\Ai'(-\zeta)/\sqrt\zeta=r\sin\Omega$, the Airy equation gives
$\Omega'=-\sqrt\zeta-\sin(2\Omega)/(4\zeta)$, so $\Omega$ is decreasing whenever
$\cos\Omega=0$, and the $k$-th zero $a_k$ of $\Ai$ has $\Omega(|a_k|)=\pi/2-k\pi$. By DLMF
9.9.6 and 9.9.18, $\frac23|a_k|^{3/2}=(4k-1)\pi/4+O(1/k)$, so $-\phi=\pi/2-k\pi+o(1)$ there,
and $n_0=0$.
\end{proof}

\begin{lemma}[Turning sum]\label{lem:turningsum}
$T_K=\sum_{m_L<m<m_R}(\pi-\varphi_m)\to\frac23K^{3/2}$ as $\lambda\to\infty$ at fixed $K$.
\end{lemma}

\begin{proof}
Blocks with $x_m\ge0$ have $\varphi_m=\pi$. For $-K<x_m<0$,
$\pi-\varphi_m=2\arccos(1-u_m/2)=2\sqrt{u_m}(1+O(u_m))=\ell^{-1}\sqrt{|x_m|}\,(1+O(K/\ell^2))$.
With grid steps $h_m=\ell^{-1}(1+o(1))$ this is a Riemann sum of $\sqrt\zeta$ on
$[0,|x_{m_R}|]$, and $|x_{m_R}|\to K$.
\end{proof}

\subsection{The near field}

\begin{proposition}{Near field}{nearfield}
Assume \Cref{hyp:Hk}. If $\kappa'=0$ let $M=\infty$; if $\kappa'\ne0$ let
$M=M(\lambda)$ satisfy $N_M/N_\lambda\to\infty$ and $N_M^{\alpha-1}/\lambda\to0$. Then
\[
  \chi_M=\sum_{m<M}\varphi_m-\frac{\kappa'\pi}{2\alpha}+o(1)\qquad(\lambda\to+\infty),
\]
where for $M=\infty$ the left side is $\psi(\lambda)$. The $o(1)$ is uniform for $\kappa'$ in
compact subsets of $(-(\alpha-1)/2,(\alpha-1)/2)$ when the other data are fixed.
\end{proposition}

\begin{proof}
This is \eqref{eq:budget} made precise. Apply \Cref{prop:adiabatic} from $m_1=m_R$ to $M$
with the split
\[
  \bar\varepsilon_m:=-2\kappa'/N_m,\qquad \bar\sigma_m:=s(\mu_m)-s(\mu_me^{-\bar\ell_m}),
\]
which keeps the drift and the smooth part of the frame change. By \Cref{lem:nearsums}
the hypotheses of \Cref{prop:adiabatic} hold for $\lambda\ge\lambda_0(K)$, and the error
$\sum_{m_R\le m<M}\mathcal E_m$ together with the boundary terms $|y_{m_R}-\chi_{m_R}|$ and
$|y_M-\chi_M|$ is at most $CK^{-3/2}+o(1)$.

\step{assembly} Blocks $m\le m_L$ have $\varphi_m=\pi$, including the finitely many small
$m$ once $\lambda>2\max_{m<m_0}A_m$. So $\sum_{m<m_R}\varphi_m=m_R\pi-T_K$. By \Cref{blem:window},
\Cref{lem:turningsum} and \Cref{prop:adiabatic},
\[
  \limsup_{\lambda\to\infty}\Bigl|\chi_M-\sum_{m<M}\varphi_m-\sum_{m_R\le m<M}\bar\varepsilon_m\tan\tfrac{\varphi_m}2\Bigr|
  \le\bigl|\Omega(K)+\tfrac23K^{3/2}-\tfrac\pi4\bigr|+C_0\eta_A(K)+CK^{-3/2},
\]
and the right side tends to $0$ as $K\to\infty$ by \Cref{lem:airyphase}.

\step{the drift} $\tan(\varphi_m/2)=(\mu_m/2)(1-\mu_m^2/4)^{-1/2}$ and
$\mu_m/2=(N_\lambda/N_m)^\alpha(1+o(1))$ on the oscillatory side. The drift sum is therefore a
Riemann sum with step $2$ in $N$ of
\[
  -\kappa'\int_{N_\lambda}^\infty\frac{(N_\lambda/N)^\alpha}{\sqrt{1-(N_\lambda/N)^{2\alpha}}}\frac{dN}{N}
  =-\frac{\kappa'}{\alpha}\int_0^1\frac{dw}{\sqrt{1-w^2}}=-\frac{\kappa'\pi}{2\alpha}
\]
($w=(N_\lambda/N)^\alpha$, $dw/w=-\alpha\,dN/N$). The integrand has an integrable singularity at
$N_\lambda$; the part below $N_{m_R}$ is $O((\ell K/N_\lambda)^{1/2})$ and the part above $N_M$ is
$O((N_\lambda/N_M)^\alpha)$.

\step{uniformity} Every constant above, and in \Cref{lem:nearsums}, depends on $\kappa'$ only
through $|\kappa'|$ and through the bound $|\tilde\varepsilon_m|\le(C+C|\kappa'|)N_m^{-1-\delta}$,
so it is bounded for $\kappa'$ in compact sets.
\end{proof}
